% Unit 041 terjemahan Bahasa Indonesia dari chapter3.tex baris 1710--1881.
% Sumber: Methods of Algebra, Volume 2, commit
% 9a5803ff77dd3257484cb177f851a73770a59dd3 (CC BY 4.0).
% stable-unit-id: o014.aljabr2.chapter3.double-complex-cohomology
% Cakupan: Bagian 3.10 lengkap.

% segment-id: o014.li2.chapter3.double-complex-cohomology.h001
\section{Kohomologi bikompleks}\label{sec:double-cplx-coh}
\hypertarget{o014.aljabr2.chapter3.double-complex-cohomology}{ }

% segment-id: o014.li2.chapter3.double-complex-cohomology.p001
Bagian ini bertujuan mengkaji kohomologi bikompleks $X$ pada arah horizontal
atau vertikal, serta hubungannya dengan kohomologi kompleks total. Hasil-hasil
terkait akan digunakan untuk mengkaji funktor turunan dari bifunktor; cara
pembuktiannya mengikuti \cite[\S 12.5]{KS06}.

% segment-id: o014.li2.chapter3.double-complex-cohomology.p002
Misalkan $\mathcal{A}$ kategori abelian. Seperti biasa, semua hasil mempunyai
perumuman yang bersesuaian ketika $\mathcal{A}$ merupakan kategori abelian
$\Bbbk$-linear.

% segment-id: o014.li2.chapter3.double-complex-cohomology.p003
Ingat bahwa Definisi~\ref{def:bicplx-F1} (atau diagram yang lebih eksplisit,
\eqref{eqn:bicplx-diagram}) memperkenalkan sepasang isomorfisme antarkategori
aditif:
% segment-id: o014.li2.chapter3.double-complex-cohomology.g001
\[\begin{tikzcd}
	\cate{C}^2(\mathcal{A}) \arrow[r, shift left, "{F_{\mathrm{I}}}"] \arrow[r, shift right, "{F_{\mathrm{II}}}"'] & \cate{C}(\cate{C}(\mathcal{A})) ,
\end{tikzcd}\]
Dengan funktor-funktor tersebut, untuk setiap bikompleks $X$ definisikan
(di bawah ini $p, q, n \in \Z$):
% segment-id: o014.li2.chapter3.double-complex-cohomology.q001
\begin{align*}
	\Hm_{\mathrm{I}}^p(X) & := (\Hm^p \circ F_{\mathrm{I}})(X)
	= \text{kompleks}\; \left[ \begin{tikzcd}[row sep=small]
		\vdots \\
		\Hm^p(X^{\bullet, q+1}, \dhori) \arrow[u, "{\dvert}" inner sep=0.6em] \\
		\Hm^p(X^{\bullet, q}, \dhori) \arrow[u, "{\dvert}" inner sep=0.6em] \\
		\vdots \arrow[u, "{\dvert}" inner sep=0.6em]
	\end{tikzcd}\right] \quad \text{(lihat Proposisi~\ref{prop:abelian-cat-cplx})}, \\
	\Hm_{\mathrm{II}}^q(X) & := (\Hm^q \circ F_{\mathrm{II}})(X) \\
	& = \text{kompleks}\; \left[ \cdots \xrightarrow{\dhori} \Hm^q(X^{p, \bullet}, \dvert) \xrightarrow{\dhori} \Hm^q(X^{p+1, \bullet}, \dvert) \xrightarrow{\dhori} \cdots \right] , \\
	\tau_{\mathrm{I}}^{\leq n} & := (F_{\mathrm{I}})^{-1} \circ \tau^{\leq n} \circ F_{\mathrm{I}}, \\
	\tau_{\mathrm{II}}^{\leq n} & := (F_{\mathrm{II}})^{-1} \circ \tau^{\leq n} \circ F_{\mathrm{II}},
\end{align*} \index[sym1]{HmIpX@$\Hm_{\mathrm{I}}^p(X)$, $\Hm_{\mathrm{II}}^q(X)$} \index[sym1]{tauI@$\tau_{\mathrm{I}}^{\leq n}$, $\tau_{\mathrm{II}}^{\leq n}, \ldots$}%
\footnote{Catatan penerjemah (O014-C054): pada kedua definisi pertama,
sumber hanya mencetak funktor komposit $\Hm^p\circ F_{\mathrm{I}}$ dan
$\Hm^q\circ F_{\mathrm{II}}$ pada ruas kanan, padahal ruas kiri merupakan
nilainya pada bikompleks $X$. Target memulihkan penerapan pada $X$.}
Selain itu, terdapat $\tilde{\tau}_{\mathrm{I}}^{\leq n}$,
$\tau_{\mathrm{I}}^{\geq n}$, $\tilde{\tau}_{\mathrm{I}}^{\geq n}$,
serta $\tilde{\tau}_{\mathrm{II}}^{\leq n}$,
$\tau_{\mathrm{II}}^{\geq n}$, dan
$\tilde{\tau}_{\mathrm{II}}^{\geq n}$. Dengan demikian, tetap terdapat
morfisme-morfisme antarfunktor
% segment-id: o014.li2.chapter3.double-complex-cohomology.q002
\begin{gather*}
	\tau_{\mathrm{I}}^{\leq n} \to \tilde{\tau}_{\mathrm{I}}^{\leq n} \to \identity_{\cate{C}^2(\mathcal{A})} \to \tilde{\tau}_{\mathrm{I}}^{\geq n} \to \tau_{\mathrm{I}}^{\geq n}, \\
	\tau_{\mathrm{II}}^{\leq n} \to \tilde{\tau}_{\mathrm{II}}^{\leq n} \to \identity_{\cate{C}^2(\mathcal{A})} \to \tilde{\tau}_{\mathrm{II}}^{\geq n} \to \tau_{\mathrm{II}}^{\geq n}.
\end{gather*}

% segment-id: o014.li2.chapter3.double-complex-cohomology.de001
\begin{definisi}\label{def:Cf-Supp}
	\index[sym1]{SuppX@$\Supp(X)$}
	\index[sym1]{C2f@$\cate{C}^2_f(\mathcal{A})$}
	Untuk bikompleks $X$, tuliskan
	$\Supp(X) := \left\{ (p,q) \in \Z^2 : X^{p,q} \neq 0 \right\}$.
	Definisikan $\cate{C}^2_f(\mathcal{A})$ sebagai subkategori penuh yang
	terdiri atas objek-objek $X$ dari $\cate{C}^2(\mathcal{A})$ yang memenuhi
	syarat berikut: untuk setiap $n \in \Z$, himpunan
	$\{ (p, q) \in \Supp(X): p + q = n \}$ berhingga.
\end{definisi}

% segment-id: o014.li2.chapter3.double-complex-cohomology.ex001
\begin{contoh}
	Jika $\Supp(X) \subset \Z_{\geq 0}^2$ (kuadran pertama), atau
	$\Supp(X) \subset \Z_{\leq 0}^2$ (kuadran ketiga), atau jika $\Supp(X)$
	merupakan gabungan berhingga dari kolom atau baris, maka $X$ merupakan
	objek dari $\cate{C}^2_f(\mathcal{A})$.
\end{contoh}

% segment-id: o014.li2.chapter3.double-complex-cohomology.p004
Kompleks total $\tot_{\oplus}(X)$ dan $\tot_{\Pi}(X)$ terdefinisi untuk
setiap objek $X$ dari $\cate{C}^2_f(\mathcal{A})$, tanpa syarat tambahan
terhadap $\mathcal{A}$. Selain itu, dalam keadaan ini
$\tot_{\oplus}(X) = \tot_{\Pi}(X)$. Dengan demikian, diperoleh funktor aditif
$\tot: \cate{C}^2_f(\mathcal{A}) \to \cate{C}(\mathcal{A})$.

% segment-id: o014.li2.chapter3.double-complex-cohomology.p005
Sekarang buat pengamatan sederhana tetapi perlu berikut. Untuk objek $X$ dari
$\cate{C}^2_f(\mathcal{A})$ dan $n \in \Z$, terdapat himpunan bagian
berhingga $S \subset \Z^2$ sedemikian sehingga
$\Hm^n\left( \tot(X) \right)$ sepenuhnya ditentukan oleh data
$\left( X^{p,q}, \dhori^{p, q} , \dvert^{p,q} \right)_{(p, q) \in S}$.

% segment-id: o014.li2.chapter3.double-complex-cohomology.le001
\begin{lema}\label{prop:tot-exact}
	Funktor kompleks total
	$\tot: \cate{C}^2_f(\mathcal{A}) \to \cate{C}(\mathcal{A})$ merupakan
	funktor eksak (Definisi~\ref{def:exact-functor}).
\end{lema}
% segment-id: o014.li2.chapter3.double-complex-cohomology.pf001
\begin{bukti}
	Misalkan $X \xrightarrow{f} Y \xrightarrow{g} Z$ merupakan barisan eksak dalam
	$\cate{C}^2_f(\mathcal{A})$. Dengan mengingat definisi-definisi terkait,
	misalnya Proposisi~\ref{prop:abelian-cat-cplx}, persoalannya ialah
	membuktikan bahwa, untuk setiap $n \in \Z$, barisan
% segment-id: o014.li2.chapter3.double-complex-cohomology.q003
	\[ \bigoplus_{p+q=n} X^{p,q} \xrightarrow{\bigoplus_{p+q=n} f^{p,q}} \bigoplus_{p+q=n} Y^{p,q} \xrightarrow{\bigoplus_{p+q=n} g^{p,q}} \bigoplus_{p+q=n} Z^{p,q} \]
	eksak dalam $\mathcal{A}$. Karena jumlah langsung tersebut berhingga,
	semuanya tereduksi pada keeksakan
	$X^{p,q} \xrightarrow{f^{p,q}} Y^{p,q} \xrightarrow{g^{p,q}} Z^{p,q}$.
\end{bukti}

% segment-id: o014.li2.chapter3.double-complex-cohomology.p006
Walaupun lema berikut tampak rumit, lema ini semata-mata merupakan pengolahan
definisi dan berlaku untuk setiap kategori aditif $\mathcal{A}$.

% segment-id: o014.li2.chapter3.double-complex-cohomology.le002
\begin{lema}\label{prop:double-cplx-tot-aux0}
	Misalkan $h \in \Z$ dan $Y$ objek dari $\cate{C}(\mathcal{A})$. Dari data
	ini, bentuk
% segment-id: o014.li2.chapter3.double-complex-cohomology.l001
	\begin{compactitem}
% segment-id: o014.li2.chapter3.double-complex-cohomology.i001
		\item objek $Y[-h]$ dari $\cate{C}(\mathcal{A})$;
% segment-id: o014.li2.chapter3.double-complex-cohomology.i002
		\item objek $Y_{\mathrm{I}}[-h]$ dari
		$\cate{C}(\cate{C}(\mathcal{A}))$: komponen berderajat $h$ ialah $Y$,
		sedangkan komponen lainnya ialah $0$.
	\end{compactitem}
	Maka terdapat isomorfisme-isomorfisme kanonik berikut dalam
	$\cate{C}(\mathcal{A})$:
% segment-id: o014.li2.chapter3.double-complex-cohomology.q004
	\begin{gather*}
		(\tot \circ F_{\mathrm{I}}^{-1}) \Cone\left( \identity_{Y_{\mathrm{I}}[-h]} \right) \simeq \Cone\left( \identity_{Y[-h]} \right), \\
		(\tot \circ F_{\mathrm{I}}^{-1}) \left( Y_{\mathrm{I}}[-h] \right) \simeq Y[-h].
	\end{gather*}
\end{lema}
% segment-id: o014.li2.chapter3.double-complex-cohomology.pf002
\begin{bukti}
	Menurut definisi, objek
	\[
		\Cone\left( \identity_{Y_{\mathrm{I}}[-h]} \right)
		\quad\text{berbentuk}\quad
		Y \xrightarrow{\identity_Y} Y
	\]
	dari kategori
	\[
		\cate{C}^{[h-1, h]}(\cate{C}(\mathcal{A})),
	\]
	yang terkonsentrasi pada derajat $h-1$ dan $h$. Oleh karena itu, dengan menetapkan
	\[
		Z := F_{\mathrm{I}}^{-1}
		\Cone\left( \identity_{Y_{\mathrm{I}}[-h]} \right),
	\]
	diperoleh bikompleks pada diagram berikut:
% segment-id: o014.li2.chapter3.double-complex-cohomology.g002
	\[\begin{tikzcd}
		n \in \Z & \vdots & \vdots \\
		(q = n-h+1) & Y^{n-h+1} \arrow[u, "d_Y"] \arrow[r, "\identity"] & Y^{n-h+1} \arrow[u, "d_Y"] \\
		(q = n-h) & Y^{n-h} \arrow[r, "\identity"] \arrow[u, "d_Y"] & Y^{n-h} \arrow[u, "d_Y"] \\
		& (p = h-1) \arrow[phantom, u, "\cdots" description, sloped] & (p = h) \arrow[phantom, u, "\cdots" description, sloped]
	\end{tikzcd}\]
	Kolom-kolom yang bersesuaian dengan $p \notin \{h-1, h\}$ semuanya
	bernilai $0$. Berdasarkan Definisi~\ref{def:total-cplx} tentang kompleks
	total, langsung diperoleh
% segment-id: o014.li2.chapter3.double-complex-cohomology.q005
	\begin{align*}
		\tot(Z)^n & = Y^{n-h+1} \oplus Y^{n-h} = Y[-h]^{n+1} \oplus Y[-h]^n \\
		& = \Cone\left( \identity_{Y[-h]} \right)^n , \\
		d_{\tot(Z)}^n & = \begin{pmatrix}
			(-1)^{h-1} d_Y^{n-h+1} & 0 \\
			\identity_{Y^{n-h+1}} & (-1)^{h} d_Y^{n-h}
		\end{pmatrix}
		= \begin{pmatrix}
			- d_{Y[-h]}^{n+1} & 0 \\
			\identity_{Y[-h]^{n+1}} & d_{Y[-h]}^n
		\end{pmatrix} \\
		& = d_{\Cone\left( \identity_{Y[-h]} \right)}^n .
	\end{align*}
	Isomorfisme kedua dapat diverifikasi dengan cara serupa, tetapi lebih
	mudah, sehingga diserahkan kepada pembaca.
\end{bukti}

% segment-id: o014.li2.chapter3.double-complex-cohomology.le003
\begin{lema}\label{prop:double-cplx-tot-aux}
	Misalkan $X$ objek dari $\cate{C}^2_f(\mathcal{A})$ dan $q \in \Z$.
% segment-id: o014.li2.chapter3.double-complex-cohomology.l002
	\begin{enumerate}[(i)]
% segment-id: o014.li2.chapter3.double-complex-cohomology.i003
		\item Morfisme alami
		$\tot\left(\tau_{\mathrm{I}}^{\leq q} X\right) \to
		\tot\left(\tilde{\tau}_{\mathrm{I}}^{\leq q} X\right)$ merupakan
		kuasi-isomorfisme;
% segment-id: o014.li2.chapter3.double-complex-cohomology.i004
		\item Terdapat barisan eksak pendek kanonik
% segment-id: o014.li2.chapter3.double-complex-cohomology.q006
		\[ 0 \to \tot\left( \tilde{\tau}_{\mathrm{I}}^{\leq q-1} (X) \right) \to \tot\left( \tau_{\mathrm{I}}^{\leq q}(X) \right) \to \Hm_{\mathrm{I}}^q(X)[-q] \to 0. \]
	\end{enumerate}
\end{lema}
% segment-id: o014.li2.chapter3.double-complex-cohomology.pf003
\begin{bukti}
	Terapkan Lema~\ref{prop:truncation-ses} dalam
	$\cate{C}(\cate{C}(\mathcal{A}))$ untuk memperoleh barisan-barisan eksak
	pendek
% segment-id: o014.li2.chapter3.double-complex-cohomology.q007
	\begin{gather*}
		0 \to \tau^{\leq q} F_{\mathrm{I}} (X) \to \tilde{\tau}^{\leq q} F_{\mathrm{I}} (X) \to \Cone\left( \identity_{\Image\left( d^q_{F_{\mathrm{I}} X} \right)_{\mathrm{I}} [-q-1]} \right) \to 0, \\
		0 \to \tilde{\tau}^{\leq q-1} F_{\mathrm{I}}(X) \to \tau^{\leq q} F_{\mathrm{I}}(X) \to \Hm^q_{\mathrm{I}}(X)_{\mathrm{I}}[-q] \to 0.
	\end{gather*}
	Arti notasi $(\cdots)_{\mathrm{I}}[\cdots]$ sama seperti dalam
	Lema~\ref{prop:double-cplx-tot-aux0}. Sekarang terapkan
	$\tot \circ F_{\mathrm{I}}^{-1}$ pada barisan-barisan eksak pendek
	tersebut. Lema~\ref{prop:tot-exact} menunjukkan bahwa $\tot$ eksak, dan
	$F_{\mathrm{I}}^{-1}$ tentu juga eksak. Dengan menggunakan
	Lema~\ref{prop:double-cplx-tot-aux0}, diperoleh barisan-barisan eksak pendek
	berikut dalam $\cate{C}(\mathcal{A})$:
% segment-id: o014.li2.chapter3.double-complex-cohomology.q008
	\begin{gather*}
		0 \to \tot\left(\tau_{\mathrm{I}}^{\leq q} X\right) \to \tot\left( \tilde{\tau}_{\mathrm{I}}^{\leq q} X \right) \to \Cone\left( \identity_{\Image( d_{F_{\mathrm{I}} X}^q )[-q-1]} \right) \to 0, \\
		0 \to \tot\left( \tilde{\tau}_{\mathrm{I}}^{\leq q-1} X \right) \to \tot\left( \tau_{\mathrm{I}}^{\leq q}X \right) \to \Hm_{\mathrm{I}}^q(X)[-q] \to 0.
	\end{gather*}
	Barisan kedua tepat merupakan (ii). Sementara itu, (i) diperoleh dari
	barisan pertama bersama Proposisi~\ref{prop:long-exact-sequence-ses} dan
	Proposisi~\ref{prop:cone-homotopy} (i).
\end{bukti}

% segment-id: o014.li2.chapter3.double-complex-cohomology.p007
Untuk setiap objek $X$ dari $\cate{C}^2(\mathcal{A})$, notasikan
bikompleks-bikompleks berikut dengan $\Hm_{\mathrm{I}}(X)$ dan
$\Hm_{\mathrm{II}}(X)$ ($p,q \in \Z$):
\index[sym1]{HmIX@$\Hm_{\mathrm{I}}(X)$, $\Hm_{\mathrm{II}}(X)$}
% segment-id: o014.li2.chapter3.double-complex-cohomology.q009
\begin{equation}\label{eqn:HIHII}\begin{aligned}
	\left( \Hm_{\mathrm{I}}(X)^{p, \bullet}, \dvert_{\Hm_{\mathrm{I}}(X)}^{p, \bullet} \right) & := \Hm_{\mathrm{I}}^p(X), &
	\dhori_{\Hm_{\mathrm{I}}(X)}^{\bullet, \bullet} := 0, \\
	\left( \Hm_{\mathrm{II}}(X)^{\bullet, q}, \dhori_{\Hm_{\mathrm{II}}(X)}^{\bullet, q} \right) & := \Hm_{\mathrm{II}}^q(X), &
	\dvert_{\Hm_{\mathrm{II}}(X)}^{\bullet, \bullet} := 0,
\end{aligned}\end{equation}
Dengan kata lain, $\Hm_{\mathrm{I}}(X)$ (atau $\Hm_{\mathrm{II}}(X)$)
diperoleh dengan mengambil kohomologi $X$ pada arah $\dhori$ (atau
$\dvert$). Keduanya merupakan funktor dari $\cate{C}^2(\mathcal{A})$ ke
dirinya sendiri. Selanjutnya, dapat didefinisikan
$\Hm_{\mathrm{II}} \Hm_{\mathrm{I}}(X)$ dan
$\Hm_{\mathrm{I}} \Hm_{\mathrm{II}} (X)$; keduanya memenuhi
$\dhori = 0 = \dvert$.

% segment-id: o014.li2.chapter3.double-complex-cohomology.th001
\begin{teorema}\label{prop:double-cplx-tot}
	Misalkan $f: X \to Y$ suatu morfisme dalam $\cate{C}^2_f(\mathcal{A})$. Jika
	morfisme yang diinduksinya
	$\Hm_{\mathrm{II}} \Hm_{\mathrm{I}} (X) \to
	\Hm_{\mathrm{II}} \Hm_{\mathrm{I}} (Y)$ merupakan isomorfisme, maka
	$\tot(f): \tot(X) \to \tot(Y)$ merupakan kuasi-isomorfisme.
	
	Kesimpulan yang sama berlaku apabila, sebagai gantinya, disyaratkan bahwa
	morfisme yang diinduksi
	$\Hm_{\mathrm{I}} \Hm_{\mathrm{II}} (X) \to
	\Hm_{\mathrm{I}} \Hm_{\mathrm{II}} (Y)$ merupakan isomorfisme.
\end{teorema}
% segment-id: o014.li2.chapter3.double-complex-cohomology.pf004
\begin{bukti}
	Syarat tersebut ekuivalen dengan mengatakan bahwa morfisme
	$\Hm_{\mathrm{I}}^n(f): \Hm_{\mathrm{I}}^n(X) \to
	\Hm_{\mathrm{I}}^n(Y)$ dalam $\cate{C}(\mathcal{A})$ merupakan
	kuasi-isomorfisme untuk setiap $n \in \Z$. Langkah pertama pembuktian
	ialah mereduksi ke kasus ketika
% segment-id: o014.li2.chapter3.double-complex-cohomology.q010
	\begin{equation}\label{eqn:double-cplx-tot-aux}
		q \ll 0 \implies \tilde{\tau}_{\mathrm{I}}^{\leq q}(X) = 0 = \tilde{\tau}_{\mathrm{I}}^{\leq q}(Y).
	\end{equation}
	Memang, misalkan $r \in \Z$. Pemenggalan dengan funktor-funktor dalam
	Lema~\ref{prop:truncation-ses} menunjukkan bahwa syarat semula juga
	mengakibatkan
	$\Hm_{\mathrm{I}}^n\left( \tau_{\mathrm{I}}^{\geq r} f \right):
	\Hm_{\mathrm{I}}^n \left(\tau_{\mathrm{I}}^{\geq r} X \right) \to
	\Hm_{\mathrm{I}}^n\left( \tau_{\mathrm{I}}^{\geq r} Y \right)$
	merupakan kuasi-isomorfisme untuk setiap $n$. Akan tetapi, untuk $n$ yang
	diberikan, definisi $\cate{C}^2_f(\mathcal{A})$ memperlihatkan bahwa ketika
	$r \ll 0$ terdapat diagram komutatif
% segment-id: o014.li2.chapter3.double-complex-cohomology.g003
	\begin{equation}\label{eqn:double-cplx-tot-aux2}\begin{tikzcd}
		\Hm^n(\tot(X)) \arrow[r, "{\Hm^n(\tot(f))}"] \arrow[d, "\sim"' sloped] & \Hm^n(\tot(Y)) \arrow[d, "\sim" sloped] \\
		\Hm^n\left( \tot(\tau_{\mathrm{I}}^{\geq r} X) \right) \arrow[r, "{\Hm^n\left( \tot (\tau_{\mathrm{I}}^{\geq r} f) \right)}"' inner sep=0.8em] & \Hm^n\left( \tot(\tau_{\mathrm{I}}^{\geq r} Y) \right)
	\end{tikzcd}\end{equation}
	Bikompleks-bikompleks pada baris kedua memenuhi
	\eqref{eqn:double-cplx-tot-aux}. Karena itu, mulai sekarang kita dapat
	mengganti $f$ dengan $\tau_{\mathrm{I}}^{\geq r} f$ dan mengasumsikan
	bahwa \eqref{eqn:double-cplx-tot-aux} berlaku. Tujuan kita ialah
	membuktikan bahwa $\Hm^n(\tot(f))$ merupakan isomorfisme untuk $n$ yang
	telah dipilih.
	
	Untuk setiap $q \in \Z$, Lema~\ref{prop:double-cplx-tot-aux} (ii)
	memberikan diagram komutatif dengan baris-baris eksak
% segment-id: o014.li2.chapter3.double-complex-cohomology.g004
	\[\begin{tikzcd}
		0 \arrow[r] & \tot\left( \tilde{\tau}_{\mathrm{I}}^{\leq q-1} X \right) \arrow[r] \arrow[d, "{\alpha_{q-1}}"] & \tot\left( \tau_{\mathrm{I}}^{\leq q} X \right) \arrow[d, "{\beta_q}"] \arrow[r] & \Hm_{\mathrm{I}}^q(X)[-q] \arrow[d, "{\Hm_{\mathrm{I}}^q(f)[-q]}" description] \arrow[r] & 0 \\
		0 \arrow[r] & \tot\left( \tilde{\tau}_{\mathrm{I}}^{\leq q-1} Y \right) \arrow[r] & \tot\left( \tau_{\mathrm{I}}^{\leq q} Y \right) \arrow[r] & \Hm_{\mathrm{I}}^q(Y)[-q] \arrow[r] & 0
	\end{tikzcd}\]%
	\footnote{Catatan penerjemah (O014-C055): sumber hanya memberi label
	$\simeq$ pada panah vertikal kanan. Hipotesis teorema menyatakan bahwa
	morfisme itu ialah kuasi-isomorfisme $\Hm_{\mathrm{I}}^q(f)[-q]$, bukan
	harus isomorfisme kompleks; target menamai morfismenya dan menyatakan sifat
	yang tepat.}
	Panah vertikal kanan merupakan kuasi-isomorfisme, dan semua panah
	vertikalnya berasal dari $f$. Kita akan membuktikan bahwa
	$\alpha_q$ dan $\beta_q$ keduanya merupakan kuasi-isomorfisme. Jika
	$\alpha_{q-1}$ merupakan kuasi-isomorfisme, maka funktorialitas barisan
	eksak panjang terkait (Proposisi~\ref{prop:long-exact-sequence-ses}),
	bersama Proposisi~\ref{prop:5-lemma}, mengakibatkan bahwa $\beta_q$ juga
	merupakan kuasi-isomorfisme. Di sisi lain,
	Lema~\ref{prop:double-cplx-tot-aux} (i) memberikan diagram komutatif
	berikut, dengan kedua panah horizontal berupa kuasi-isomorfisme:
% segment-id: o014.li2.chapter3.double-complex-cohomology.g005
	\[\begin{tikzcd}
		\tot\left( \tau_{\mathrm{I}}^{\leq q} X \right) \arrow[r] \arrow[d, "{\beta_q}"'] & \tot\left( \tilde{\tau}_{\mathrm{I}}^{\leq q} X \right) \arrow[d, "{\alpha_q}"] \\
		\tot\left( \tau_{\mathrm{I}}^{\leq q} Y \right) \arrow[r] & \tot\left( \tilde{\tau}_{\mathrm{I}}^{\leq q} Y \right)
	\end{tikzcd}\]
	Oleh karena itu, jika $\beta_q$ merupakan kuasi-isomorfisme, maka
	$\alpha_q$ juga demikian. Karena untuk $q \ll 0$, $\alpha_q$ tidak lain
	daripada $0 \rightiso 0$, secara rekursif diperoleh bahwa $\alpha_q$ dan
	$\beta_q$ merupakan kuasi-isomorfisme untuk setiap $q$.
	
	Ingat bahwa $n \in \Z$ telah dipilih. Dengan alasan yang serupa dengan
	\eqref{eqn:double-cplx-tot-aux2}, ketika $q \gg 0$ terdapat diagram
	komutatif
% segment-id: o014.li2.chapter3.double-complex-cohomology.g006
	\[\begin{tikzcd}
		\Hm^n(\tot(X)) \arrow[r, "{\Hm^n(\tot(f))}"] & \Hm^n(\tot(Y)) \\
		\Hm^n\left( \tot(\tau_{\mathrm{I}}^{\leq q} X) \right) \arrow[r, "{\Hm^n(\beta_q)}"' inner sep=0.8em, "\sim"] \arrow[u, "\sim" sloped] & \Hm^n\left( \tot(\tau_{\mathrm{I}}^{\leq q} Y) \right) \arrow[u, "\sim"' sloped]
	\end{tikzcd}\]
	Jadi, $\Hm^n(\tot(f))$ memang merupakan isomorfisme.
	
	Terakhir, argumen yang sama berlaku untuk kasus
	$\Hm_{\mathrm{I}} \Hm_{\mathrm{II}} (X) \rightiso
	\Hm_{\mathrm{I}} \Hm_{\mathrm{II}} (Y)$. Sudut pandang lainnya ialah
	menukar $\Hm_{\mathrm{I}}$ dan $\Hm_{\mathrm{II}}$ dengan menggunakan
	Proposisi~\ref{prop:double-cplx-swap}.
\end{bukti}

% segment-id: o014.li2.chapter3.double-complex-cohomology.p008
Pembuktian di atas cukup panjang. Contoh~\sourcecrossref{eg:double-cplx-tot-ss}{pada \S~5.6} akan
memberikan pembuktian lain yang didasarkan pada barisan spektral.

% segment-id: o014.li2.chapter3.double-complex-cohomology.co001
\begin{korolari}\label{prop:acyclic-assembly}
	Misalkan $X$ objek dari $\cate{C}_f^2(\mathcal{A})$. Jika $X$ eksak pada
	baris, yakni jika $\left( X^{\bullet, q}, \dhori \right)$ eksak untuk
	setiap $q \in \Z$, maka $\tot(X)$ eksak. Demikian pula, jika $X$ eksak
	pada kolom, maka $\tot(X)$ eksak.
\end{korolari}
% segment-id: o014.li2.chapter3.double-complex-cohomology.pf005
\begin{bukti}
	Jika $X$ eksak pada baris, maka
	$\Hm_{\mathrm{II}} \Hm_{\mathrm{I}}(X) = 0$. Karena itu,
	Teorema~\ref{prop:double-cplx-tot} mengakibatkan bahwa
	$\tot(X) \to \tot(0) = 0$ merupakan kuasi-isomorfisme; dengan kata lain,
	$\tot(X)$ eksak.

	Argumen untuk kasus eksak pada kolom sama persis. Sebagai alternatif,
	kedua kasus dapat dialihkan satu ke yang lain dengan menggunakan
	Proposisi~\ref{prop:double-cplx-swap}.
\end{bukti}
